\begin{figure}[ht]
  \centering
  \begin{tikzpicture}
    \begin{groupplot}[group style={group size=2 by 1, horizontal sep=1.8cm},
        tesi, width=.46\textwidth, height=5.8cm, xmin=0, xmax=110000,
        scaled x ticks=false, xtick={10018,50000,100000}, xticklabels={10.018,50.000,100.000},
        x tick label style={font=\footnotesize, rotate=30, anchor=north east},
        xlabel={documenti}, ymin=0]
      \nextgroupplot[title={\footnotesize Memoria a riposo (MB)}, legend to name=leg:crescita, legend columns=2]
        \addplot[color=black, mark=square*, mark size=1.6pt] table[x=documenti, y=mem_es] {dati/crescita.dat};
        \addplot[color=black, mark=triangle*, mark size=2pt, dashed] table[x=documenti, y=mem_k] {dati/crescita.dat};
        \legend{Elasticsearch, Koskidex}
      \nextgroupplot[title={\footnotesize Indice su disco (MB)}]
        \addplot[color=black, mark=square*, mark size=1.6pt] table[x=documenti, y=disco_es] {dati/crescita.dat};
        \addplot[color=black, mark=triangle*, mark size=2pt, dashed] table[x=documenti, y=disco_k] {dati/crescita.dat};
    \end{groupplot}
  \end{tikzpicture}

  \ref{leg:crescita}
  \caption[Memoria e indice su disco al crescere dell'archivio]{Memoria a riposo e indice su disco al crescere dell'archivio, con gli
  stessi atti copiati fino a 50.000 e 100.000 documenti. Elasticsearch ha
  l'heap fissato a 1~GB, quindi la sua memoria non cresce; quella di Koskidex sì.
  Dati: \texttt{latex/dati/crescita.dat}.}
  \label{fig:crescita}
\end{figure}
